\section{\tool{}}\label{sec:design}

Section~\ref{sec:study} shows that binding flaws are decided by comparing two
artifacts that no single analyzer reads together. \tool{} mechanizes that
comparison with an LLM \emph{agent} that must build the two ledgers of
Section~\ref{sec:def-ledgers} explicitly, a tool-less \emph{post-stage} that
keeps the agent from losing what it found, and an \emph{execution loop} that
turns warnings into attack tests. Its central design rule is
\textbf{additivity}: a stage may add warnings and evidence but never remove or
alter an earlier stage's, which is checked mechanically at every stage
boundary. Additivity makes every stage separately measurable and every gain
attributable (Figure~\ref{fig:pipeline}).

\begin{figure*}[t]
\centering
\begin{tikzpicture}[
  st/.style={draw=blue!60!black, rounded corners=2pt, fill=blue!5, minimum height=12mm, text width=19mm, align=center,
             font=\sffamily\scriptsize, inner sep=2pt},
  pst/.style={st, draw=green!50!black, fill=green!7},
  lst/.style={st, draw=gray, fill=white},
  arr/.style={-{Latex[length=1.6mm]}, gray!70!black, thick}, node distance=3.2mm]
\node[lst] (s0) {\textbf{Inventory}\\[1pt]entries $E_i$ with effects; statement producers $S_j$};
\node[st, right=of s0] (s1) {\textbf{Effect ledger}\\[1pt](entry, effect, value): bound / unbound};
\node[st, right=of s1] (s2) {\textbf{Statement ledger}\\[1pt]public value: what fixes it?};
\node[st, right=of s2] (s3) {\textbf{Follow-through, attacker pass}\\[1pt]final warnings};
\node[pst, right=5mm of s3] (s4) {\textbf{No-drop}\\[1pt]re-adjudicate discarded suspicions};
\node[pst, right=of s4] (s5) {\textbf{Multi-location}\\[1pt]one finding, every root-cause site};
\node[lst, right=5mm of s5] (s6) {\textbf{Execution loop}\\[1pt]E1 claim, E2 mock test, E3 real proof};
\foreach \a/\b in {s0/s1,s1/s2,s2/s3,s3/s4,s4/s5,s5/s6} \draw[arr] (\a) -- (\b);
\draw[arr, dashed] (s2.south) to[bend left=28] node[below, font=\sffamily\tiny]{coverage re-prompt} (s1.south);
\begin{scope}[on background layer]
\node[draw=blue!60!black, dashed, rounded corners=3pt, fit=(s0)(s3), inner sep=4pt,
  label={[font=\sffamily\scriptsize, blue!60!black]above:Stage A: one hermetic agent session (read-only copy, capped)}] {};
\node[draw=green!50!black, dashed, rounded corners=3pt, fit=(s4)(s5), inner sep=4pt,
  label={[font=\sffamily\scriptsize, green!40!black]above:post-stage: tool-less}] {};
\end{scope}
\end{tikzpicture}
\caption{\tool{} pipeline. Stage A drives the agent through coverage-checked
ledgers. The post-stage sees Stage A's ledgers and the tail of its
transcript, not the repository. The loop attaches executed attack tests.
Every stage only adds; new warnings carry a \texttt{finding\_id}, a source
and a location role (proof side, consumer, helper).}
\label{fig:pipeline}
\end{figure*}

\subsection{Stage A: ledger-guided review}\label{sec:design-stagea}

A plain LLM asked to ``find binding bugs'' reads the consumer, notices the
most salient issue and stops. The failure is coverage: a target has many
entry points, each with several effect values, and a flaw hides in one row
of a table the model never writes down. Stage A turns the definition into
that table. One session of an LLM runs in a hermetic sandbox: a read-only
private copy of the snapshot, file-reading tools and deterministic helpers,
with no network and no write access. A controller drives it through fixed
passes and parses tagged output lines (\texttt{LEDGER}, \texttt{STMT},
\texttt{WARNING}) after every turn (Appendix~\ref{app:extra}).
\begin{enumerate}
\item \textbf{Inventory.} A deterministic helper lists the entry points that
reach a verifier and commit effects ($E_1,\dots$) and the producers of public
statements ($S_1,\dots$): verifier calls, public-input arrays, circuit
\texttt{main} signatures and guest commit sites. The list is the coverage
target.
\item \textbf{Effect ledger.} One row per (entry, effect, value) with the
value's provenance (statement, hash-committed input, storage, calldata,
sender) and a verdict, \emph{bound} or \emph{unbound}, with a reason. An
entry without a row is re-prompted, at most twice.
\item \textbf{Statement ledger.} For each public value: what the circuit or
guest constrains it to mean, which consumer value it should fix, and whether
the consumer compares the two. This pass catches P3--P5, where the consumer
is right but the statement is not.
\item \textbf{Follow-through and attacker pass.} The agent runs the playbook
of each unbound row's pattern, follows values into libraries, helpers and
record writers, takes the view of an attacker holding an observed proof, and
emits final warnings: pattern, file and function, effect, unbound value,
explanation and related locations.
\end{enumerate}
Cost guards skip optional passes above a soft cap and force the final answer
at a hard cap. A typical effect-ledger row from a benchmark run reads
(identifiers generalized):
\begin{lstlisting}
E1 Pool.processBundle | executed calls of an op with 0 joinsplits
   value: actions of a proofless op | source: op in bundle | UNBOUND
   why: validateOperation does not require a joinsplit, so an op
   needs no proof if another op in the bundle has one
\end{lstlisting}

\subsection{How agents lose what they find}\label{sec:design-lose}

Reading Stage A's ledgers against its final answers on development data
showed two ways in which a flaw the agent had recorded fails to reach the
output.

\para{Placement}
A binding flaw has two sides, and the agent reports it on one. Given a guest
that commits a prover-chosen starting state and a contract that stores the
result, it reports at the contract; the statement ledger already records the
guest-side gap, but the final answer folds it into the contract warning. When
the root cause is in the guest, the report points the developer to the wrong
fix and, under strict scoring, counts as a miss.

\para{Non-reporting}
The agent writes an unbound row or drafts a warning and then lets it go: it
withdraws it with a reason that sounds prudent but is invalid under the
threat model (the consumer ``is not in this repository'', the value at stake
seems too small, the behavior is intended design, an exclusion that does not
apply), or never carries the row into the answer. In an audit, where
a false negative costs far more than a false positive, this is the wrong
prior.

\subsection{Post-stage: no-drop and multi-location}\label{sec:design-post}

\para{No-drop}
The first step collects three kinds of candidates from Stage A: drafts
emitted during the session but absent from the final answer
($W\setminus W_f$), effect rows with verdict \emph{unbound}, and statement
rows marked unconstrained. One tool-less LLM call receives the candidates, the
final warnings, the closing part of the transcript, the exclusion list and a
closed list of \textbf{admissible drop reasons}. Writing $\ell(\cdot)$ for a
primary location (file, function) and $\nu(\cdot)$ for the unbound value, a
candidate $c$ may be dropped only if
\[
\begin{aligned}
\textsf{drop}(c)\iff{}&\exists w\in W_f:\ \nu(w)=\nu(c)\wedge \ell(w)=\ell(c)\\
&\vee\ \textsf{chk}(c)\ \vee\ \textsf{excl}(c)\ \vee\ \textsf{noact}(c),
\end{aligned}
\]
where \textsf{chk} requires a file:line check that binds the value on every
path (or shows that the attack only makes verification fail), \textsf{excl} a
numbered exclusion that genuinely applies, and \textsf{noact} that the
attacker changes nothing while keeping a valid proof. Any other reason,
including every dismissal of Section~\ref{sec:design-lose}, restores $c$ as a
warning where the value enters the effect path (Appendix~\ref{app:nodrop}).
The same-location clause turns placement into recall: ``already part of
another finding'' is inadmissible when the locations differ. The step sees no
repository, so it cannot invent locations; it can only overrule the agent's
own decisions.

\para{Multi-location}
The second step groups warnings into \emph{findings} (same unbound value,
same missing check) and lists, per finding, up to three root-cause locations
tagged by role (proof side, consumer, helper). Each location is validated
against a function definition in the snapshot and emitted as a sibling
warning sharing the finding's \texttt{finding\_id}. Volume, precision and
false alarms are therefore counted in findings.

\subsection{Execution loop and evidence}\label{sec:design-loop}

A warning is a claim; auditors write a proof of concept. For each unbound
row with an on-chain entry point, the loop drafts a Foundry scaffold once per
(contract, entry), checks that an honest call succeeds, and then submits the
same accepted proof twice from the same state: once honestly and once with
the attacker's change. The test succeeds when both calls are accepted and the
committed effects differ in recipient, amount, state written or count, which
is condition C3 made executable; a success is re-run with the honest value
restored and must then show no difference. Attempt~0 comes from the pattern's
template without the LLM (resubmit from a second account, change a context
argument, replay, add the field modulus to the replay key, swap the verifier);
up to three further attempts are written by the agent from the labeled
outcomes of earlier ones. Evidence has three levels: \emph{E1} (claim),
\emph{E2} (a test against a mock verifier that accepts exactly the honest
public inputs succeeded) and \emph{E3} (a real proof from the project's own
prover, against the unmodified verifier and consumer, succeeded). Six oracle
guards close loopholes we found by running the loop on patched code and on
24 adversarial harnesses written to fool it, such as an owner-only call that
looks like a replay because the scaffold deployed from the honest account
(Appendix~\ref{app:guards}); with them, none of the 10 rows attacked at a
fixed or hard-negative location and none of the 24 harnesses reaches E2. The loop never removes a warning, and
whether a test succeeded is computed from the two executions, never read from
the agent's summary. Circuit-only and guest-only flaws without an on-chain
attack call stay at E1.

\subsection{Extensions we built and rejected}\label{sec:design-ext}

Four further steps target the misses that remain. All are generic, sit
behind their own flags, and were measured under the same rule as the kept
steps (Section~\ref{sec:eval-anatomy}). \emph{Knowledge} loads a pitfall
catalogue (34 flaw mechanisms, 17 look-alikes) and stack-specific files for
the detected proof system (RISC~Zero, SP1, Noir, Circom, gnark, halo2,
verifier integration, e-mail proofs, light clients, TEE attestation) into
Stage A. \emph{Challenge} runs a second tool-using session that re-opens every
non-unbound verdict on an effect value under a list of forbidden dismissals.
\emph{Inventory} enumerates every private input of a zkVM guest and asks
whether it can steer a committed output. \emph{Probes} ask for domain
separation in every proven message, parser differentials on authenticated
fields, and anchor checks on every library getter.

\subsection{Implementation and isolation}\label{sec:design-impl}

\tool{} is about 6,700 lines of Python. The agent and the post-stage use
\texttt{claude-sonnet-5-5} at high reasoning effort through a command-line
client run non-interactively with a tool allowlist and write and network
access denied; any tool-using model can be substituted. Helpers parse
Solidity, Circom, Noir and Rust with the Solidity compiler's AST, tree-sitter
and regular expressions. A Stage A run is capped at 20 minutes. A
\emph{limit guard} treats any model call that returns in under 10\,s at zero
cost, with an error or with an empty answer as a failed run that is never
cached or scored: provider usage limits interrupted our study three times,
and before the guard existed 52 failed calls had been cached as empty
results, which would have silently lowered the recall of the affected
configurations. Every version was built by a separate process that could read
only development data, and an automated gate rejects project names, file
names, addresses and benchmark identifiers in prompts and code.
